\chapter{Appendix A: Muscle Force and Joint Torque Calculations}\label{app:muscle_calcs}
\graphicspath{{../Figures/92-AppendixA/}{figs/92-AppendixA/}}

This appendix documents the reference implementation of the force, torque, and compliance equations used in Chapters~\ref{ch:muscles} and~\ref{ch:torque}. When present in MATLAB code, file names are given so that any number in this dissertation can be traced to the code that produced it.

\section{Human muscle force calculations}
\subsection{OpenSim Thelen2003 and Millard2012 Formulations}

The Thelen2003 and Millard2012 equilibrium muscle models describe active and passive fiber force, pennation, and a series tendon \citep{thelen_adjustment_2003,millard_flexing_2013}. Their common force balance is
\begin{equation}
 F_T(l_T)=F_M(l_M,\dot l_M,a)\cos\alpha,
 \qquad l_{MT}=l_T+l_M\cos\alpha,
 \label{eq:app_muscle_equilibrium}
\end{equation}
where $a$ is activation, $\alpha$ is pennation, and $l_M$, $l_T$, and $l_{MT}$ are fiber, tendon, and musculotendon lengths. The cosine projects fiber force onto the tendon direction. When tendon compliance is enabled, fiber length is a state determined by the equilibrium calculation rather than the full musculotendon path length.

In the damped Millard formulation, fiber force has the form
\begin{equation}
 F_M=F_0\left[a f_L(\tilde l_M)f_V(\tilde v_M)
 + f_{PE}(\tilde l_M)+\beta\tilde v_M\right],
 \label{eq:app_millard_force}
\end{equation}
with normalized fiber length $\tilde l_M$, normalized velocity $\tilde v_M$, maximum isometric force $F_0$, and damping coefficient $\beta$. The active, passive, and tendon curves and the treatment of singular configurations differ between the Thelen and Millard models. The OpenSim implementation can also disable tendon compliance or activation dynamics, so the model class alone does not establish which dynamics were exercised in an analysis \citep{millard_flexing_2013}.

\subsubsection{Thelen activation and contraction dynamics}

The following seven relations correspond, in order, to Equations~(1)--(7) of \citet{thelen_adjustment_2003}. They specify the activation state, passive and active force--length curves, tendon force, and fiber velocity. Bars denote normalization by maximum isometric fiber force $F_0^M$ or optimal fiber length $L_0^M$, as appropriate.

Activation follows
\begin{equation}
 \frac{da}{dt}=\frac{u-a}{\tau_a(a,u)},
 \label{eq:app_thelen1}
\end{equation}
with an activation-dependent time constant
\begin{equation}
 \tau_a(a,u)=\begin{cases}
 \tau_{\mathrm{act}}(0.5+1.5a),&u>a,\\
 \tau_{\mathrm{deact}}/(0.5+1.5a),&u\leq a.
 \end{cases}
 \label{eq:app_thelen2}
\end{equation}
Here $u\in[0,1]$ is neural excitation, $a$ is muscle activation, and $\tau_{\mathrm{act}}$ and $\tau_{\mathrm{deact}}$ set the activation and deactivation time scales.

The normalized passive fiber force and active force--length multiplier are
\begin{equation}
 \bar F^{PE}=\frac{\exp\!\left[k^{PE}(\bar L^M-1)/\varepsilon_0^M\right]-1}
 {\exp(k^{PE})-1},
 \label{eq:app_thelen3}
\end{equation}
\begin{equation}
 f_l=\exp\!\left[-\frac{(\bar L^M-1)^2}{\gamma}\right].
 \label{eq:app_thelen4}
\end{equation}
The tendon relation is
\begin{equation}
 \bar F^T=\begin{cases}
 \displaystyle\frac{\bar F^T_{\mathrm{toe}}}{\exp(k_{\mathrm{toe}})-1}
 \left[\exp\!\left(\frac{k_{\mathrm{toe}}\varepsilon^T}{\varepsilon^T_{\mathrm{toe}}}\right)-1\right],
 &\varepsilon^T\leq\varepsilon^T_{\mathrm{toe}},\\[6pt]
 k_{\mathrm{lin}}(\varepsilon^T-\varepsilon^T_{\mathrm{toe}})+\bar F^T_{\mathrm{toe}},
 &\varepsilon^T>\varepsilon^T_{\mathrm{toe}}.
 \end{cases}
 \label{eq:app_thelen5}
\end{equation}
The fiber-velocity relation and its shortening/lengthening denominator are
\begin{equation}
 V^M=(0.25+0.75a)V^M_{\max}\frac{\bar F^M-a f_l}{b},
 \label{eq:app_thelen6}
\end{equation}
\begin{equation}
 b=\begin{cases}
 a f_l+\bar F^M/A_f,&\bar F^M\leq a f_l,\\[4pt]
 \displaystyle\frac{(2+2/A_f)(a f_l\bar F^M_{\mathrm{len}}-\bar F^M)}
 {\bar F^M_{\mathrm{len}}-1},&\bar F^M>a f_l.
 \end{cases}
 \label{eq:app_thelen7}
\end{equation}

Here $\bar F^M$ is the normalized active fiber force entering the force--velocity relation, $V^M=\dot L^M$ is negative during shortening, and $V^M_{\max}$ is the maximum shortening-speed scale. Table~\ref{tab:app_thelen_symbols} defines the curve parameters. A shape factor is a dimensionless parameter controlling a curve's curvature or width; it is not a stiffness in N/m. For example, increasing $\gamma$ broadens the active force--length curve while preserving its peak at $\bar L^M=1$. Changing $k^{PE}$ changes the passive curve's shape while retaining unit force at strain $\varepsilon_0^M$.

\begin{table}[htbp]
\centering\small
\caption{Length, strain, and shape parameters in the Thelen formulation.}
\label{tab:app_thelen_symbols}
\begin{tabularx}{\textwidth}{@{}lX@{}}
\toprule
Symbol & Definition \\
\midrule
$\bar L^M=L^M/L_0^M$ & Fiber length relative to optimal fiber length. \\
$\varepsilon_0^M$ & Passive fiber strain at force $F_0^M$. \\
$k^{PE}$ & Passive force--length exponential shape factor (5 in the paper). \\
$\gamma$ & Active force--length width parameter (0.45 in the paper). \\
$\varepsilon^T$ & Tendon strain relative to tendon slack length. \\
$\varepsilon_0^T$ & Tendon strain at force $F_0^M$. \\
$\varepsilon^T_{\mathrm{toe}}$ & Strain at the end of the nonlinear tendon region; $0.609\varepsilon_0^T$. \\
$\bar F^T_{\mathrm{toe}}$ & Normalized force at that transition; 0.33. \\
$k_{\mathrm{toe}},\ k_{\mathrm{lin}}$ & Tendon exponential shape factor (3) and linear slope ($1.712/\varepsilon_0^T$). \\
$A_f$ & Force--velocity shape factor (0.25 in the paper). \\
$\bar F^M_{\mathrm{len}}$ & Maximum normalized lengthening-force factor. \\
\bottomrule
\end{tabularx}
\end{table}

The passive expressions describe tensile loading; they do not imply that a slack muscle or tendon transmits compression. These are the original paper's constitutive relations. Numerical implementations must also handle slack states, minimum activation, pennation limits, and the inversion of the force--velocity relation. Consequently, the original equations alone do not specify every numerical safeguard in an OpenSim implementation.

\subsubsection{Millard equilibrium and damping}

The equilibrium relation of \citet{millard_flexing_2013} can be written in normalized form as
\begin{equation}
 \left[a f_L(\tilde l_M)f_V(\tilde v_M)+f_{PE}(\tilde l_M)\right]\cos\alpha
 -f_T(\tilde l_T)=0.
 \label{eq:app_millard_equilibrium}
\end{equation}
The damped formulation adds $\beta\tilde v_M$ inside the brackets, as in Equation~\ref{eq:app_millard_force}. The functions $f_L$, $f_V$, $f_{PE}$, and $f_T$ describe active force--length, force--velocity, passive fiber force, and tendon force. Millard's smooth characteristic curves and damping address numerical behavior while retaining the series force balance. Figure~3 of that paper illustrates the equilibrium, damped-equilibrium, and rigid-tendon formulations. A rigid-tendon model removes tendon stretch as a state; it does not simply set the passive fiber force to zero.


\subsection{OpenSim Moment Arms and Human Torque Targets}
\label{sec:app_opensim_targets}

The human torque targets of Section~\ref{sec:optimization} are OpenSim
tendon-excursion moment arms: for each muscle, the moment arm about the
knee is the partial derivative of muscle-tendon length with respect to
the knee coordinate, with positive tensile force and torque positive in the coordinate direction, $r(\theta_{k}) = -\partial l_{\mathrm{MT}} /
\partial\theta_{k}$, evaluated across the range of motion. The exported
analysis tables archived with the test data
(\nolinkurl{OpenSim_Vasti_MomentArm.txt}, \nolinkurl{OpenSim_Vasti_Results.txt},
\nolinkurl{OpenSim_Bifem_MomentArm.txt}, \nolinkurl{OpenSim_Bifem_Results.txt})
carry the per-muscle moment arms and the corresponding
maximum-isometric torque curves (moment arm multiplied by each
muscle's maximum isometric force) for vastus intermedius,
lateralis, and medialis, and for the biceps femoris short head; the
optimization loads the vastus medialis and biceps femoris short-head
columns as its targets (\texttt{buildKneeExtContext20mm.m} lines
119--140; \nolinkurl{Knee_Flexor_Data_20mm.m} line 546). For an unconstrained planar rotation, this
derivative definition and the geometric perpendicular-foot construction
of Equation~(\ref{eq:app_perpfoot}) below are equivalent for a straight
muscle path crossing a planar revolute joint; the constrained and moving-path-point geometry is evaluated by OpenSim rather than by substituting the robot-side straight-path construction.

\section{Converted MuJoCo Muscle Formulation}
\label{sec:muscle-models}

\subsection{MuJoCo Muscle Actuator}

The converted model uses MuJoCo muscle actuators, which combine activation-dependent active force with passive tensile force. Let $L$ denote the routed musculotendon length, $L_T$ the constant biological-tendon length, and $L_0$ the optimal fiber length inferred from the actual and normalized operating ranges. The normalized muscle length is $\bar L=(L-L_T)/L_0$. The signed actuator force is
\begin{equation}
 F_{\mathrm{act}}=-F_0\left[a f_L(\bar L)f_V(\dot L/(L_0v_{\max}))+f_P(\bar L)\right].
 \label{eq:mj-muscle}
\end{equation}
Here $v_{\max}$ is the maximum shortening speed in optimal fiber lengths per second. The negative sign denotes a pulling actuator force; both active and passive terms contribute to tension. The rigid tendon has a constant length, but no elastic state. MuJoCo also omits pennation, so its differences from the OpenSim equilibrium formulations are not limited to tendon compliance \citep{mujoco_muscle_docs_237}.

\section{BPA Force Model Implementation}

\subsection{Strain Measures}
The musculotendon path length $l_{LMT}$ is the sum of straight segments between route via points, transformed into a common frame at each frame crossing. The contractile strain used by the force model subtracts the artificial tendon and both end-cap fittings from the path length before normalization by the resting length,
\begin{equation}
    \varepsilon \;=\; \frac{l_{0} - \left(l_{LMT} - l_{ten} - 2\,l_{fitting} - \chi_{0}\right)}{l_{0}},
    \label{eq:app_strain}
\end{equation}
where $\chi_{0}$ is the constant length offset of Section~\ref{sec:improved_model} (zero in the uncorrected model). [Provenance: \texttt{MonoPamDataExplicit.m} get.Contraction, lines 210--217; \texttt{minimizeFlxPin.m} lines 415--426. End-cap fitting length $l_{fitting}=\qty{25.4}{\mm}$ is a constructor default, \texttt{MonoPamDataExplicit.m} line 69.]

The maximum free-load contraction at \qty{620}{\kPa} is $\varepsilon_{620} = (l_{0}-l_{620})/l_{0}$, where $l_{620}$ is the BPA length at \qty{620}{\kPa} with no external load (code variable \texttt{KMAX}; \texttt{MonoPamDataExplicit.m} line 278). The relative strain is $\varepsilon^{*} = \varepsilon / \varepsilon_{620}$ (line 281).

\subsection{Normalized Force Surface}
The normalized isometric force surface of Equation~(\ref{eq:Fstar}) is implemented in closed form in \nolinkurl{Functions/f_festo.m} (lines 1--26),
\begin{equation}
    F^{*} = c_{0}\left(e^{-c_{1}\varepsilon^{*}} - 1\right)
            + P^{*}\,e^{-c_{2}(\varepsilon^{*})^{2}},
    \qquad F^{*}\le 0 \;\Rightarrow\; 0,
    \label{eq:app_fstar}
\end{equation}
with $P^{*} = P/\qty{620}{\kPa}$. The zero-clamp $F^{*}\le 0 \Rightarrow 0$ is applied by the callers (e.g.\ \nolinkurl{Functions/festo4.m} line 35), not inside \nolinkurl{f_festo.m} itself. The coefficients were fitted with \texttt{fittype} nonlinear least squares and least-absolute-residual robustness (\nolinkurl{Testing_Data/bpaFits.m}, lines 45--60 and 115--140, lower bounds $[0,0,0]$) and are reported in Table~\ref{tab:Fstar}. The code additionally holds a 40\,mm fit not used in this dissertation ($c_0 = 0.1224$, $c_1 = 10.47$, $c_2 = 2.023$; \nolinkurl{f_festo.m} lines 17--20).

At runtime the torque model may alternatively evaluate a spline interpolation of the manufacturer's datasheet curves (\nolinkurl{Functions/festo4.m}, backed by \texttt{FestoLookup.mat}), with force clamped to zero for $\varepsilon^{*}>1$ or $F^{*}<0$. Equation~(\ref{eq:app_fstar}) replaced this lookup in the corrected model.

\subsection{Maximum Force}
The resting-length-dependent maximum force of Equation~(\ref{eq:Fmax_3D}) is implemented in \nolinkurl{Functions/maxBPAforce.m} (10\,mm: line 46; 20\,mm: line 57; duplicate for 20\,mm in \nolinkurl{Functions/Fmax20.m}) with the fixed \qty{7.5}{\mm} end-contact offset. The fitted 20\,mm coefficients carry confidence intervals $a_1 \in [1.4745, 1.5011]$\,\si{\N\per\kPa} and $a_2 \in [0.0238, 0.0258]$\,\si{\per\kPa\per\metre} (\texttt{Fmax20.m} lines 26--27). Total force is the product $F = F^{*}(\varepsilon^{*},P^{*})\,F_{620}(l_{0})$, resolved into a vector along the force direction (\texttt{MonoPamDataExplicit.m} lines 283--287).

\section{Muscle Path Length and Moment Arm}

Path length accumulates the Euclidean norm of each segment between via points, applying the body-frame transformation whenever the path crosses between the femur and tibia frames (\texttt{MonoPamDataExplicit.m} lines 78--95, total at lines 122--133).

The vector moment arm is implemented as the perpendicular foot of the joint origin onto the line of action,
\begin{equation}
    \mathbf{r} \;=\; \mathbf{p} \;-\; \hat{\mathbf{u}}\left(\hat{\mathbf{u}}\cdot\mathbf{p}\right),
    \label{eq:app_perpfoot}
\end{equation}
where $\mathbf{p}$ is the vector from the joint origin to the muscle-path crossing point and $\hat{\mathbf{u}}$ is the unit force direction (\texttt{MonoPamDataExplicit.m} lines 156--168; identical implementations in \texttt{MonoMuscleData.m} lines 111--123, the \texttt{Mom} function of \texttt{minimizeFlxPin.m} at line 539, and \nolinkurl{MonoPam_mult.m} lines 186--201). This is the same quantity as the cross-product projection of Equation~(\ref{eq:project_ma}); the planar scalar moment arm about the joint axis is the magnitude of the in-plane components, computed as \texttt{hypot(mA(:,1), mA(:,2))} (\texttt{predictKneeFlexor20mm.m} line 217). Torque about the joint $z$-axis is $\mathbf{m} = \mathbf{r}\times\mathbf{f}$, evaluated only where the BPA is in tension (values are discarded for $\varepsilon < -0.02$, the excessive stretching region; \texttt{MonoPamDataExplicit.m} line 308).

\section{Compliance Model Implementation}

The compliance-corrected muscle length entering Equation~(\ref{eq:Lm}) is computed as $l_{m} = l_{LMT} - \chi_{0} - l_{ten} - 2\,l_{fitting}$ (\texttt{minimizeFlxPin.m} lines 415--426; identically \texttt{minimizeExtX3.m} line 345). In the two-bracket fit of record, the insertion-bracket stiffness matrix and the projected scalar stiffness along the force direction are
\begin{equation}
    \mathbf{K}_{br} = \mathrm{diag}(\chi_{1},\,\chi_{2},\,\chi_{1}),
    \qquad
    k_{\hat{u},br} = \left(\hat{u}^{\,T}\mathbf{K}_{br}^{-1}\hat{u}\right)^{-1},
    \label{eq:app_kbr}
\end{equation}
matching Equation~(\ref{eq:bktstiffness}), and the origin-side second bracket contributes an additional scalar compliance $c_{b2}$ in series (Appendix~\ref{app:model_geometry}, Table~\ref{tab:app_stiffness_arrays}). As executed by the evaluator of record, both bracket frames are assembled pitch-only, at \texttt{minimizeFlxPin.m} line 289 (insertion) and line 306 (origin), with the two-rotation assemblies carried as commented alternatives; the insertion array is $[X_{1},X_{2},X_{1}]$ (line 387) and the origin array $[X_{1},X_{1},X_{2}]$ (line 400). The two-rotation-capable twin of the same block is \texttt{minimizeFlxPin2brk.m} lines 387--418. An earlier arm-1 array $[X_{1},X_{2},X_{2}]$ is superseded and survives only in the provenance comment at line 387, and the legacy single-bracket evaluator uses $[X_{1},X_{2},X_{1}]$ (\texttt{minimizeFlxPinX3.m} line 267).

The artificial tendon is modeled as an axial spring with stiffness
\begin{equation}
    k_{ten} \;=\; \frac{n\,A_{c}\,E}{L_{ten}},
    \label{eq:app_tendon}
\end{equation}
with $n=2$ parallel cable strands ($\times 6$ for the 20\,mm assembly), $A_{c} = \qty{1.51e-6}{\metre\squared}$ (19-strand cable), and $E = \qty{193e9}{\pascal}$ (\nolinkurl{MonoPam_mult.m}, function \texttt{Spr()}, lines 627--637; the evaluators obtain the spring rate through \nolinkurl{Spr(bpa_i)}, \texttt{minimizeFlxPin.m} line 156). Series combination with the projected bracket compliance gives $k_{eq} = \left(k_{\hat{u},br}^{-1} + k_{ten}^{-1} + c_{b2}\right)^{-1}$ (\texttt{minimizeFlxPin.m} lines 394--409).

The scalar compliance displacement $\delta$ of Equation~(\ref{eq:complianceforcebalance}) is found numerically with \texttt{fzero} on the bracket $[0,\, l_{m}-l_{620}]$,
\begin{equation}
    F_{620}(l_{0})\,F^{*}\!\left(\varepsilon^{*}(\delta),P^{*}\right) - k_{eq}\,\delta = 0,
    \label{eq:app_forcebalance}
\end{equation}
with $\delta \equiv 0$ when the uncorrected relative strain already exceeds unity (\texttt{minimizeFlxPin.m} lines 434--441). The vector bracket deflection is then $\boldsymbol{\epsilon}_{br} = \mathbf{K}_{br}^{-1}\left(\norm{\mathbf{f}}\,\hat{\mathbf{u}}\right)$ (line 461), the cable stretch is $\norm{\mathbf{f}}/k_{ten}$ (line 478), and the attachment point, path, moment arm, and torque are recomputed. For the two-BPA flexor grouping, the coupled balance shares one common force between both muscles and is solved analogously on $[0, F_{620}]$ (\nolinkurl{MonoPam_mult.m} lines 556--612).

\section{Wrapping-Loss Implementation}

For the pinned-knee extensor, Equation~(\ref{eq:deltaL}) is implemented segment-wise over the wrap arc with radii $R_1 = \qty{12}{\mm}$ and $R_2 = \qty{40}{\mm}$ and segment boundaries $\theta_{1} = \ang{27}$ and $\theta_{2} = \ang{-23}$ (\texttt{minimizeExtX3.m}, nested \texttt{Contraction}, lines 327--370). For the biomimetic-knee 20\,mm cell, the arc-length factor is the routed bend measure $b(\theta_{k})$ (arc length swept by the wrap, built with the \texttt{wRap} bend radius in \texttt{buildKneeFlexorRoute20mm.m} lines 265--268), giving $\Delta l = \chi_{3}\, b(\theta_{k})\,(1-\varepsilon^{*})^{2}$ (\nolinkurl{MonoPam_mult.m} lines 398--441), with a hardcoded-geometry fallback for degenerate routes (lines 428--434). The loss enters only the force-producing strain; the kinematic (prediction) strain excludes it (\texttt{minimizeExtX3.m} lines 236--240).

